\chapter{LITERATURE REVIEW}

\section{General Overview}
The response of a building during an earthquake depends not only on the superstructure but also on the behaviour of the foundation and supporting ground. Wolf's formulation of dynamic soil--structure interaction established the coupled nature of these components and the limitations of treating the foundation as perfectly fixed \cite{wolf1985}. More recent studies continue to show that soil flexibility can alter vibration characteristics, deformation demand, and force distribution, particularly when the supporting ground is relatively soft \cite{poudel2024,kc2025}.

SSI is commonly discussed in terms of kinematic and inertial interaction. Kinematic interaction concerns the way a foundation modifies the motion that would exist in the free field, while inertial interaction arises from forces generated by the vibrating structure and transmitted through the foundation to the soil \cite{wolf1985,kc2025}. Both mechanisms are relevant to performance assessment because they can change the response quantities used to judge serviceability and damage, including displacement and inter-storey drift.

\section{Seismic Vulnerability of Buildings in Nepal}
Nepal's seismic vulnerability reflects both regional tectonic activity and characteristics of the building stock. Gautam et al. identified recurring structural and construction deficiencies in Nepalese buildings, while studies following the 2015 Gorkha earthquake documented the consequences of these weaknesses during strong shaking \cite{gautam2016,adhikari2022}. Paudel et al. likewise examined vulnerability in Nepalese reinforced-concrete residential buildings and highlighted the sensitivity of seismic performance to structural configuration \cite{paudel2024seismic}.

Local ground conditions add another source of uncertainty. Deep alluvial deposits, particularly in the Kathmandu Valley, can modify earthquake motion and contribute to differences between fixed-base predictions and the response of structures supported on deformable soil \cite{poudel2024,adhikari2022}. Poudel and Chaulagain's nonlinear investigation and the simplified study by Timilsina et al. both indicate that soil flexibility can materially affect displacement and drift demand in reinforced-concrete buildings \cite{poudel2024,timilsina2023}.

\section{Hospital Buildings as Critical Facilities}
The importance of hospital seismic performance in Nepal has been recognized for more than two decades. Guragain and Dixit's assessment of Nepalese hospitals showed that many facilities required significant improvement to achieve acceptable earthquake performance \cite{guragain2004}. More recent work by Subedi compared seismic demand in an RC hospital building using IS 1893:2016 and NBC 105:2020, illustrating how design provisions can influence calculated response \cite{subedi2022}.

These studies are especially relevant because the acceptable outcome for a hospital is not simply avoidance of collapse. The building must retain sufficient structural and functional capacity to support emergency services after an earthquake. This requirement makes deformation-sensitive measures, foundation response, and the distribution of structural demand important alongside conventional force-based checks \cite{guragain2004,subedi2022}.

\section{Soil--Structure Interaction Theory}
Wolf describes SSI as a dynamic interaction problem in which the motion of the foundation and the response of the superstructure influence each other through the supporting medium \cite{wolf1985}. On stiff ground, the fixed-base idealization may provide a reasonable approximation; as soil compliance increases, foundation translation and rotation become part of the total structural response. KC et al. summarize these effects within a broader review of ground-motion modification by soil, structures, and topography \cite{kc2025}.

The practical consequence is a change in effective stiffness, period, damping, and deformation pattern. Santisi d'Avila and Lopez-Caballero, for example, examined nonlinear interaction in a coupled soil--structure system, while Adhikari et al. demonstrated the importance of site condition in fragility analysis of RC buildings on alluvial deposits \cite{santisi2016,adhikari2022}. These findings support the use of flexible-base models when the study objective is to compare the response of the same building across different soil conditions.

\section{Foundation Spring and Flexible-Base Modelling}
SSI may be represented either through direct modelling of the soil domain or through a substructure approach in which foundation impedances are introduced into the structural model \cite{wolf1985,gajan2008}. The substructure approach is attractive for building analysis because translational and rotational stiffnesses can be defined at the foundation level while the superstructure remains within a conventional structural-analysis program.

The classical Winkler concept idealizes the supporting soil as a set of independent linear springs for which the local reaction is proportional to the local foundation displacement. In its simplest vertical form this is written as $p=k_s w$, where $p$ is contact pressure, $w$ is local displacement, and $k_s$ is the modulus of subgrade reaction. Contemporary SSI reviews describe the Winkler model as a simple spring representation of the soil domain, while also noting its limitations for reproducing the continuum response of soil \cite{krishnan2024}. The present study uses this spring-idealization concept at the structural-model level, but it does not obtain the six foundation stiffness components from a single subgrade-reaction modulus.

Instead, the directional spring magnitudes are obtained from the elastic-half-space foundation-impedance expressions assembled by Gazetas. Gazetas provides algebraic expressions and charts for surface and embedded rigid foundations, including vertical, horizontal, rocking, and torsional modes, as functions of soil properties and foundation geometry \cite{gazetas1991}. Thus, the modelling adopted here is most precisely described as a Winkler-type discrete spring representation whose translational, rocking, and torsional stiffness values are derived from the Gazetas elastic-half-space impedance formulation. Gajan et al. and Raychowdhury extend related spring concepts to nonlinear shallow-foundation mechanisms such as uplift, gapping, sliding, and cyclic settlement \cite{gajan2008,raychowdhury2008}; those nonlinear foundation mechanisms are outside the present linear spring idealization.

\section{Influence of Soil Condition on Seismic Response}
The effect of SSI is not uniformly beneficial or adverse. Period lengthening can shift a structure to a different region of the response spectrum, which may reduce some force measures while increasing displacement or drift. Timilsina et al. reported noticeable changes in seismic response when soil flexibility was introduced using a simplified method, and Krishnan and Sivakumar similarly emphasized the influence of SSI on the stability and response of RC structures \cite{timilsina2023,krishnan2024}.

Poudel and Chaulagain investigated nonlinear SSI effects in irregular RC buildings and found that the interaction between soil and structure can substantially modify response compared with a fixed-base representation \cite{poudel2024}. Taken together, these studies indicate that the direction and magnitude of SSI effects should be established from the response quantities of the specific model rather than assumed in advance.

\section{Performance-Based Seismic Assessment}
Performance-based assessment evaluates seismic demand in relation to defined performance objectives rather than relying solely on elastic force checks. FEMA 356 provides a widely used framework for performance levels such as Immediate Occupancy, Life Safety, and Collapse Prevention \cite{fema356}. The approach places particular emphasis on deformation and component behaviour, which makes it compatible with an SSI study where foundation flexibility can alter drift and displacement.

Janous and El Ghoulbzouri examined the vulnerability of irregular RC buildings while considering SSI, and more recent studies have continued to incorporate soil effects into performance-oriented evaluation \cite{janous2024,oggu2025}. For hospitals, these methods are relevant because the expected post-earthquake condition is tied to continued service, not only to collapse resistance \cite{subedi2022}.

\section{Pushover Analysis and Performance Levels}
Nonlinear static pushover analysis estimates lateral capacity by incrementally increasing lateral load and tracking the resulting roof displacement, base shear, and nonlinear response. FEMA~356 defines a nonlinear static procedure in which the structural model is pushed to a target displacement and the resulting forces and deformations are checked against the applicable performance criteria \cite{fema356}. The method therefore provides a direct basis for relating the capacity response to a displacement-based performance objective. Ahmadi and Celebi specifically examined the influence of SSI on pushover response in reinforced-concrete frames, showing that support flexibility can affect the inferred performance \cite{ahmadi2018}.

Performance levels are commonly interpreted through deformation and component acceptance criteria. Within the present study, pushover results are used to compare the capacity response of the four support idealizations, while the reported performance coordinates are interpreted within the FEMA~356 target-displacement framework. The comparison focuses on how the base shear and roof displacement associated with the reported performance state vary between fixed, hard-soil, medium-soil, and soft-soil support conditions.

\section{Fragility Analysis and Damage States}
Fragility curves express the probability that a defined damage state will be exceeded at a specified intensity measure. Lognormal models are commonly used to relate intensity measures such as peak ground acceleration or spectral acceleration to damage-state exceedance probability \cite{giordano2021,adhikari2022}. The choice of engineering demand parameter and damage thresholds therefore has a direct influence on the resulting fragility relationship.

For Nepalese buildings, Giordano et al. developed empirical fragility models for school buildings, while Adhikari et al. analysed RC buildings with brick infills in a high-seismic region with alluvial deposits \cite{giordano2021,adhikari2022}. Kircil and Ethemoglu and Oggu and Gopikrishna examined the effect of SSI on damage probability and fragility in multistorey RC frames \cite{kircil2025,oggu2025}. These studies provide the context for interpreting the drift-based fragility trends developed from the nonlinear time-history results in this thesis.

\section{Research Gap}
The literature establishes that SSI can influence the seismic response of reinforced-concrete buildings, but Nepal-focused work on mid-size hospital buildings remains limited relative to studies of residential, school, or generic RC frames \cite{adhikari2022,paudel2024seismic,giordano2021,yadav2023}. Existing hospital studies emphasize vulnerability and code-based demand, whereas fewer studies integrate flexible-base modelling with response-spectrum demand, foundation response, member-force redistribution, nonlinear pushover capacity, and fragility within the same hospital prototype \cite{guragain2004,subedi2022}.

The present thesis addresses this gap by comparing a G+5 hospital model under fixed, hard-soil, medium-soil, and soft-soil support conditions. The comparison is carried through modal and response-spectrum response, foundation and member demands, nonlinear static capacity, and a preliminary fragility assessment using nonlinear time-history results from twelve ground motions. This integrated structure allows the effect of support flexibility to be examined consistently across several engineering demand parameters.

\section{Summary of Literature Review}
The reviewed studies show that SSI changes the dynamic system rather than producing a single predictable change in seismic demand. Period, displacement, drift, shear, foundation response, and member demand can respond differently as soil flexibility increases. This is especially important for hospitals, where deformation and post-earthquake functionality are central to performance assessment.

The literature also supports the use of pushover and fragility methods to move from elastic response toward performance interpretation. For the present study, these findings justify a model-by-model comparison of the four support conditions and provide the theoretical basis for interpreting the linear, nonlinear, and preliminary fragility results presented in the following chapters.
